% TOP_PROBLEM: {"id": 2598, "title": "Kourovka Notebook Problem 21.89", "category_id": 17, "category": {"id": 17, "name": "group_theory", "display_name": "Group Theory", "description": "Problems about groups, group actions, representations, and related algebraic structures.", "slug": "group-theory", "order_index": 17, "created_at": "2026-07-31T15:26:25.671Z"}, "status": "open", "problem_number": "KOU-21.89", "set_id": 10}
% FIRST_POSTED: 2026-10-01
% ENTRY_KIND: statement_only
% UnsolvedMath ID 2598; source catalogue code KOU-21.89
% !TeX program = lualatex
% Definition and sources only; this file is not an LLM solution attempt.
\documentclass[11pt,a4paper]{article}
\usepackage[margin=25mm]{geometry}
\usepackage{fontspec}
\setmainfont{lmroman10-regular.otf}[BoldFont=lmroman10-bold.otf,
 ItalicFont=lmroman10-italic.otf,BoldItalicFont=lmroman10-bolditalic.otf]
\usepackage{amsmath,amssymb,mathtools}
\usepackage{unicode-math}
\setmathfont{latinmodern-math.otf}
\AtBeginDocument{\let\setminus\smallsetminus}
\usepackage{xurl,hyperref}
\hypersetup{colorlinks=true,urlcolor=blue}
\setcounter{secnumdepth}{0}
\setlength{\parindent}{0pt}
\setlength{\parskip}{5pt}
\setlength{\emergencystretch}{3em}
\begin{document}
\section*{Kourovka Notebook Problem 21.89}\label{2598}
{\small UnsolvedMath ID 2598 · KOU-21.89 · Group Theory · Kourovka Notebook - New Problems, Issue 21}
\subsection{Definitions and mathematical statement}
For a nonnegative integer $n$, an integer partition of $n$ is a finite nonincreasing list $\lambda_1\ge\cdots\ge\lambda_r\ge1$ with $\sum_i\lambda_i=n$; for $n=0$ allow the empty list. Let $p(n)$ be the number of these lists. For positive $n$, the factorial is $n!=1\cdot2\cdots n$.

The problem asks whether
\[
 \forall n\in\mathbb Z,\quad n\ge40\ \Longrightarrow\ p(n)\nmid n!.
\]
Divisibility here is integer divisibility: $p(n)\mid n!$ would mean that $n!=a\,p(n)$ for some positive integer $a$. Thus one counterexample at any integer $n\ge40$ would refute the assertion, while checking only an initial finite interval does not settle it.

In group-theoretic terms, let $S_n$ be the group of all permutations of $n$ points. It has order $n!$. Every permutation decomposes into disjoint cycles, and the multiset of their lengths is a partition of $n$. Two permutations are conjugate exactly when these length multisets coincide, so $k(S_n)=p(n)$. The question is consequently whether the number of conjugacy classes of $S_n$ ever divides its order beyond degree $39$. This equivalence fixes both the counting convention for partitions and the strict lower threshold on $n$.
\subsection{Short English statement}
For every integer $n\ge40$, does the partition number $p(n)$ fail to divide $n!$?
\subsection{Sources}
\begin{itemize}
\item[S1] ULAM.AI and the UnsolvedMath Contributors, supplied problems.json, record 2598 (KOU-21.89), Kourovka Notebook - New Problems, Issue 21. Catalogue identity and source transcription; CC BY 4.0.
\url{https://huggingface.co/datasets/ulamai/UnsolvedMath}.
\item[S2] The Kourovka Notebook, 21st edition, version 47 (30 September 2026), new problems, Problem 21.89. Expanded formulation of the supplied catalogue statement.
\url{https://arxiv.org/pdf/1401.0300v47}.
\end{itemize}
\end{document}
